\documentclass[10pt,letterpaper]{article}

% ============================================================
% PACKAGES
% ============================================================

\usepackage[
    letterpaper,
    left=0.472in,
    right=0.472in,
    top=0.208in,
    bottom=0.208in
]{geometry}

\usepackage{fontspec}
\usepackage{xcolor}
\usepackage{hyperref}
\usepackage{enumitem}
\usepackage{ragged2e}
\usepackage{titlesec}

% ============================================================
% FONT
% ============================================================

\setmainfont{Carlito}

% ============================================================
% HYPERLINK SETTINGS
% ============================================================

\definecolor{linkblue}{HTML}{1155CC}

\hypersetup{
    colorlinks=true,
    linkcolor=linkblue,
    urlcolor=linkblue,
    citecolor=linkblue,
    pdfborder={0 0 0}
}

% ============================================================
% GENERAL FORMAT
% ============================================================

\pagestyle{empty}

\setlength{\parindent}{0pt}
\setlength{\parskip}{0pt}

% ============================================================
% SECTION FORMAT
% ============================================================

\titleformat{\section}
    {\bfseries\fontsize{10}{11}\selectfont}
    {}
    {0pt}
    {}

\titlespacing*{\section}
    {0pt}
    {6pt}
    {2.5pt}

% ============================================================
% CUSTOM BULLET
% ============================================================

\newcommand{\resumeBullet}[1]{%
    \noindent
    \hspace{4pt}
    {\fontsize{9}{10.3}\selectfont
    \textbullet\hspace{5pt}#1}
    \par
    \vspace{1pt}
}

% Simple two-column heading line (no LaTeX table -> safer for ATS text extraction)
\newcommand{\resumeHeaderLine}[2]{%
    \noindent{\bfseries #1}\hfill{\bfseries #2}\par
}

% ============================================================
% DOCUMENT
% ============================================================

\begin{document}

% ============================================================
% HEADER
% ============================================================

\begin{center}

{\bfseries\fontsize{19}{20}\selectfont
VISHNU KUMAR
}

\vspace{1.5pt}

{\itshape\fontsize{10.5}{11.5}\selectfont
Software Engineer -- Python \textbar\ Backend \& API Development
}

\vspace{1.5pt}

{\fontsize{9.3}{10.3}\selectfont
Agra, Uttar Pradesh, India
\;|\;
+91 90278 72285
\;|\;
\href{mailto:vishnu9027872285@gmail.com}
{vishnu9027872285@gmail.com}
\;|\;
\href{https://vishnukumar.pro}
{\textcolor{linkblue}{Portfolio}}
\;|\;
\href{https://github.com/HelloVishnu04}
{\textcolor{linkblue}{GitHub}}
}

\end{center}

% ============================================================
% PROFESSIONAL SUMMARY
% ============================================================

\section*{PROFESSIONAL SUMMARY}

\begin{justify}
{\fontsize{9}{10.3}\selectfont
Final-year B.Tech Computer Science student and Python software engineer with
hands-on experience designing, coding, and deploying REST APIs, backend
services, and modular software systems using Python (FastAPI, Flask, Django
REST Framework) and object-oriented design principles. Experienced in
building scalable microservices, integrating relational and NoSQL databases,
and automating data pipelines, with production deployment experience using
Docker, Git/GitHub version control, and CI/CD (GitHub Actions). IBM
SkillsBuild AI/ML Virtual Intern and Oracle Cloud Infrastructure Data Science
Professional certified, with software delivered independently, in an
internship setting, and through competitive hackathons.
}
\end{justify}

% ============================================================
% TECHNICAL SKILLS
% ============================================================

\section*{TECHNICAL SKILLS}

{\fontsize{9}{10.3}\selectfont

\textbf{Programming Languages:}
Python, JavaScript/TypeScript, SQL

\par\vspace{1pt}

\textbf{Software Engineering:}
Object-Oriented Programming, REST API Design, Microservices, Modular System
Design, Debugging \& Troubleshooting, Version Control (Git/GitHub)

\par\vspace{1pt}

\textbf{Backend \& APIs:}
FastAPI, Flask, Django REST Framework, JWT Authentication, WebSockets,
Swagger/OpenAPI Documentation

\par\vspace{1pt}

\textbf{Databases:}
PostgreSQL, MySQL, MongoDB, Redis

\par\vspace{1pt}

\textbf{Frontend:}
React, MERN Stack, HTML/CSS

\par\vspace{1pt}

\textbf{DevOps \& Cloud:}
Docker, GitHub Actions (CI/CD), AWS, Oracle Cloud Infrastructure (OCI)

\par\vspace{1pt}

\textbf{Data \& ML (applied):}
Pandas, NumPy, Scikit-learn, TensorFlow, Hugging Face, LangChain

}

% ============================================================
% PROFESSIONAL EXPERIENCE
% ============================================================

\section*{PROFESSIONAL EXPERIENCE}

\resumeHeaderLine{IBM SkillsBuild (PBEL Program) -- AI/ML Virtual Intern}{Jun 2026 -- Aug 2026}

\vspace{1.5pt}

{\fontsize{9}{10.3}\selectfont
{\itshape Project}
\href{https://github.com/HelloVishnu04/Recon.git}
{\textcolor{linkblue}{\itshape
: Recon -- Cybersecurity Threat-Intelligence Backend System
}}
}

\vspace{2pt}

\resumeBullet{
Designed and implemented a Python REST API (FastAPI) that serves a Random
Forest classification model to detect 5 network-attack categories, with
Swagger-documented endpoints for production consumption.
}

\resumeBullet{
Built modular, object-oriented data-preprocessing pipelines in Python
(Pandas, Scikit-learn) implementing feature scaling, missing-value
imputation, and correlation-based feature selection on high-dimensional
telemetry data.
}

\resumeBullet{
Engineered a dual-server backend architecture -- a FastAPI inference service
plus a Flask dashboard -- with pickle-based model versioning, delivering a
working system within a 2-month Agile sprint.
}

% ============================================================
% ENGINEERING PROJECTS
% ============================================================

\section*{ENGINEERING PROJECTS}

% ------------------------------------------------------------
% S.A.F.E.
% ------------------------------------------------------------

\noindent
\href{https://github.com/HelloVishnu04/S.A.F.E.git}
{\textcolor{linkblue}{\textbf{S.A.F.E. -- Real-Time Python Backend Pipeline}}}
{\bfseries\ (2026)}
\hspace{4pt}
{\itshape -- Python, REST API, TensorFlow}

\vspace{2pt}

\resumeBullet{
Designed and deployed a real-time, end-to-end ingestion-to-inference backend
pipeline in Python with automated model training, versioning, and a
scalable REST API for serving predictions.
}

% ------------------------------------------------------------
% SentiScan
% ------------------------------------------------------------

\noindent
\href{https://github.com/HelloVishnu04/Sentiment-Analysis.git}
{\textcolor{linkblue}{\textbf{SentiScan -- Flask REST API \& NLP Engine}}}
{\bfseries\ (2025)}
\hspace{4pt}
{\itshape -- Python, Flask, REST API}

\vspace{2pt}

\resumeBullet{
Designed and deployed a Flask-based REST microservice serving a
Bidirectional LSTM sentiment-classification model, achieving 92.4\%
accuracy across 50K real-world reviews.
}

\resumeBullet{
Built and exposed real-time inference endpoints in Python for low-latency,
production-style text-sentiment scoring.
}

% ------------------------------------------------------------
% MediPredict
% ------------------------------------------------------------

\noindent
\href{https://github.com/HelloVishnu04/MediPredict-.git}
{\textcolor{linkblue}{\textbf{MediPredict -- Python ML Recommendation Engine}}}
{\bfseries\ (2025)}
\hspace{4pt}
{\itshape -- Python, Scikit-learn, XGBoost}

\vspace{2pt}

\resumeBullet{
Engineered automated model-selection logic in Python across 7 medical
prediction domains, achieving a 91.3\% mean accuracy through ensemble
logic.
}

\resumeBullet{
Built an interactive Python (Streamlit) interface delivering real-time,
explainable predictions.
}

% ------------------------------------------------------------
% GPT LLM Playground
% ------------------------------------------------------------

\noindent
\href{https://github.com/HelloVishnu04/GPT_LLM_Playground.git}
{\textcolor{linkblue}{\textbf{GPT LLM Playground -- Transformer Built From Scratch}}}
{\bfseries\ (2026)}
\hspace{4pt}
{\itshape -- Python, TensorFlow}

\vspace{2pt}

\resumeBullet{
Implemented a GPT-1-style transformer architecture from scratch in Python,
building the attention and tokenization mechanics underlying modern
generative-AI systems.
}

\resumeBullet{
Designed modular, reusable Python pipelines for data preprocessing,
training, and inference, enabling reproducible experimentation.
}

% ============================================================
% EDUCATION
% ============================================================

\section*{EDUCATION}

\resumeHeaderLine{B.Tech in Computer Science -- Anand Engineering College, Agra (AKTU)}{2023 -- 2027 (Expected)}

\vspace{2pt}

\resumeHeaderLine{Intermediate (PCM) \& High School (Science) -- Shri Narsingh Bhagwan Inter College, Agra}{2020 -- 2022}

% ============================================================
% CERTIFICATIONS
% ============================================================

\section*{CERTIFICATIONS}

\resumeBullet{
Oracle Cloud Infrastructure (OCI) Data Science Professional -- Oracle (2025)
}

\resumeBullet{
Exploratory Data Analysis, Gold Certification (NSQF Level 5) -- NASSCOM
FutureSkills Prime, MeitY, Govt. of India (2026)
}

\resumeBullet{
SQL Certificate (Intermediate) -- HackerRank (2024)
}

% ============================================================
% ACHIEVEMENTS
% ============================================================

\section*{ACHIEVEMENTS}

\resumeBullet{
Maintains 17+ public GitHub repositories spanning Python software
projects, REST APIs, and applied machine-learning systems
(\href{https://github.com/HelloVishnu04}
{\textcolor{linkblue}{github.com/HelloVishnu04}}).
}

\resumeBullet{
2nd Place -- Hack The Box TechTrix Cybersecurity Competition; Top 5
Finalist (Pan-India) -- Virtual Lab Hackathon.
}

\resumeBullet{
Solved 200+ DSA problems across LeetCode \& HackerRank, strengthening
object-oriented problem solving, data structures, and algorithm design in
Python.
}

\end{document}
